\setcounter{section}{0}

\Needspace{5\baselineskip}
\hypertarget{basic-ideas-of-policy-based-reinforcement-learning}{%
\section{Basic Ideas of Policy-Based Reinforcement Learning}\label{basic-ideas-of-policy-based-reinforcement-learning}}

\Needspace{5\baselineskip}
\hypertarget{i.-limitations-of-value-based-methods}{%
\subsection{I. Limitations of Value-Based Methods}\label{i.-limitations-of-value-based-methods}}

Value-based reinforcement learning relies on learning Q(s,a), the value of an action in a given state. This requires enumerating all X actions to select the best, with discrete outputs. For continuous actions (such as robotic arm control), a network can approximate at most as many discrete actions as it has output nodes, producing stiff motion. Otherwise, fitting Q requires comparing expectations for almost infinitely many policies to select the optimum, requiring almost infinitely many parameters. We therefore introduce policy-based reinforcement learning.

\Needspace{5\baselineskip}
\hypertarget{ii.-policy-based-reinforcement-learning}{%
\subsection{II. Policy-Based Reinforcement Learning}\label{ii.-policy-based-reinforcement-learning}}

\Needspace{4\baselineskip}
\begin{enumerate}
\def\labelenumi{\arabic{enumi}.}
\tightlist
\item
  Overall idea
\end{enumerate}

Directly learn the probability distribution of actions in each state and sample accordingly. As in learning to ride a bicycle (or precise continuous robot movements), the best expected state value cannot be calculated; instead, let the agent continually select policies in each state (more precisely, policy probability distributions: a pathfinding puppy might learn to turn left with 70\% probability and right with 30\% at an intersection). Each complete journey yields cumulative reward G (unlike the Q function's expectation under the optimal policy). The ultimate objective is a policy π maximizing expected cumulative reward J. When policy is represented by a parameter vector θ (such as neural network weights), the problem becomes finding optimal θ.

\Needspace{4\baselineskip}
\begin{enumerate}
\def\labelenumi{\arabic{enumi}.}
\setcounter{enumi}{1}
\tightlist
\item
  Neural network function
\end{enumerate}

The network takes state s and outputs action values or action-related parameters. It fits π, the probability distribution of the action-value vector or action-related parameters given the state (a density for continuous values). It is experience-based, with experience stored in network parameters θ.

\Needspace{4\baselineskip}
\begin{enumerate}
\def\labelenumi{\arabic{enumi}.}
\setcounter{enumi}{2}
\tightlist
\item
  Objective function (loss function)
\end{enumerate}

Let the initial state be fixed at s0. The objective is to maximize expected cumulative reward J=E\_τ\textasciitilde π(R(τ)) for a complete action sequence. R(τ) is the reward received over the full sequence τ, and the expectation uses the trajectory distribution induced by π.

\Needspace{4\baselineskip}
\begin{enumerate}
\def\labelenumi{\arabic{enumi}.}
\setcounter{enumi}{3}
\tightlist
\item
  Policy gradient
\end{enumerate}

\Needspace{5\baselineskip}
The gradient of expected cumulative reward J (the expectation of R(τ)) with respect to parameter vector θ:

Gradient expression: start from the objective definition and transform it using the integral form of expectation and the log-likelihood trick.

\[
\nabla_\theta J(\theta)
=\nabla_\theta\int P(\tau|\theta)R(\tau)d\tau
=\int\nabla_\theta P(\tau|\theta)R(\tau)d\tau
\]

\Needspace{5\baselineskip}
Using \(\nabla_\theta P(\tau|\theta)=P(\tau|\theta)\nabla_\theta\log P(\tau|\theta)\) gives:

\[
\nabla_\theta J(\theta)
=\int P(\tau|\theta)R(\tau)\nabla_\theta\log P(\tau|\theta)d\tau
=\mathbb{E}_{\tau\sim\pi_\theta}\left[R(\tau)\nabla_\theta\log P(\tau|\theta)\right]
\]

We have transformed an integral of a gradient into an expectation, allowing the gradient to be approximated by sampling.

\Needspace{5\baselineskip}
Decompose action sequence τ into the ``elementary actions'' at each time step:

\Needspace{5\baselineskip}
Factorizing trajectory probability: the probability of trajectory \(\tau=(s_0,a_0,s_1,a_1,\cdots)\) is:

\[
P(\tau|\theta)=p(s_0)\prod_{t=0}^{T}\pi_\theta(a_t|s_t)P(s_{t+1}|s_t,a_t)
\]

\Needspace{5\baselineskip}
Only the policy term \(\pi_\theta(a_t|s_t)\) depends on \(\theta\). Taking the log gradient therefore eliminates environmental dynamics (initial distribution \(p(s_0)\) and transitions \(P(s_{t+1}|s_t,a_t)\)):

\[
\nabla_\theta\log P(\tau|\theta)=\sum_{t=0}^{T}\nabla_\theta\log\pi_\theta(a_t|s_t)
\]

(Here s\_i is the state and a\_i the action, corresponding to a product of (action-selection probability * state-transition probability).)

\Needspace{5\baselineskip}
Final form: substituting into the expectation yields the classic policy gradient theorem:

\[
\nabla_\theta J(\theta)
=\mathbb{E}_{\tau\sim\pi_\theta}
\left[
R(\tau)\cdot\sum_{t=0}^{T}\nabla_\theta\log\pi_\theta(a_t|s_t)
\right]
\]

What is \(\nabla_\theta\log\pi_\theta(a_t|s_t)\)?

\begin{itemize}
\tightlist
\item
  A vector indicating ``how to move in parameter space to increase the probability of choosing \(a_t\) in \(s_t\).''
\item
  Simply put: it seeks to make this action more probable.
\end{itemize}

This term determines the gradient direction, the update direction for vector θ. R determines the sign and magnitude (scaling the gradient vector): larger positive rewards encourage gradient ascent in that direction, while negative rewards do the opposite. The gradient is an expectation, so sampling under π\_θ yields an unbiased estimate.

Since the expectation is indexed by π\_θ, the objective changes with the policy, and the updated policy relates directly to the behavior policy. This is an online learning algorithm.

\FloatBarrier
\Needspace{5\baselineskip}
\hypertarget{references-54}{%
\subsection{References}\label{references-54}}

\vspace{0.25em}
\begin{itemize}
\tightlist
\item
  Williams, R. J. (1992). \href{https://doi.org/10.1007/BF00992696}{Simple statistical gradient-following algorithms for connectionist reinforcement learning}. Machine Learning, 8, 229-256.
\item
  Sutton, R. S., McAllester, D., Singh, S., \& Mansour, Y. (1999). \href{https://proceedings.neurips.cc/paper/1999/hash/464d828b85b0bed98e80ade0a5c43b0f-Abstract.html}{Policy Gradient Methods for Reinforcement Learning with Function Approximation}. NeurIPS 1999.
\item
  Sutton, R. S., \& Barto, A. G. (2018). \href{http://incompleteideas.net/book/the-book-2nd.html}{Reinforcement Learning: An Introduction}. MIT Press.
\end{itemize}

\clearpage
